\section{Operators and finite-dimensional laws}\label{sec:operators}

Fix a realization in the probability-one event of
Proposition~\ref{prop:inputs}, and retain its extracted sequence. All limits in
this section and the next are on these fixed measured graphs and this fixed
measured sphere. Let $Y^n$ be the stationary chain with generator $-G_n$, and set
$\xi^n_t=z_n(Y^n_t)$. We first identify its finite-dimensional limits. The
centered inverses contain the required positive-time information, provided that
the constants and the orthogonal complements of the discrete spaces are
treated separately.

\begin{proposition}\label{prop:operators}
For every $s>0$ and $\alpha>0$, the following limits hold in operator norm on
$H$:
\begin{align}
 J_n e^{-sG_n}J_n^*&\longrightarrow J e^{-sG}J^*,
 \label{eq:semigroup-convergence}\\
 J_n\alpha(\alpha+G_n)^{-1}J_n^*
 &\longrightarrow J\alpha(\alpha+G)^{-1}J^*.
 \label{eq:resolvent-convergence}
\end{align}
Consequently the finite-dimensional laws of $\xi^n$ converge to those of
the stationary Liouville Brownian motion $B^h$.
\end{proposition}

\begin{proof}
Write $\Pi$ for the orthogonal projection onto the constants in $H$. On the
mean-zero part of $J_nL^2(\mu_n)$, the operator $2K_n$ is the inverse of $G_n$.
On the constants and on $(J_nL^2(\mu_n))^\perp$, it is zero. Therefore
\begin{equation}\label{eq:functional-calculus}
 J_n e^{-sG_n}J_n^*=\Pi+\psi_s(2K_n),
 \qquad
 \psi_s(y)=
 \begin{cases}e^{-s/y},&y>0,\\0,&y=0.\end{cases}
\end{equation}
In particular, both sides vanish on the orthogonal complement of the discrete
space, while both act as the identity on constants. The same identity holds
for $G$ and $K$, with no extra orthogonal complement since $J$ is unitary.
Hilbert--Schmidt convergence of $K_n$ implies operator-norm convergence and
places all their spectra in a common bounded interval. Since $\psi_s$ is
continuous on that interval, continuous functional calculus proves
\eqref{eq:semigroup-convergence}. Similarly,
\[
 J_n\alpha(\alpha+G_n)^{-1}J_n^*
   =\Pi+r_\alpha(2K_n),
 \qquad r_\alpha(y)=\frac{\alpha y}{1+\alpha y},
\]
and the corresponding identity for $G$ proves
\eqref{eq:resolvent-convergence}. Neither argument asserts operator-norm
convergence at time zero.

To identify finite-dimensional laws, take continuous real functions
$f_0,\ldots,f_k$ on $\Sph$ and times
$0\le t_0<t_1<\cdots<t_k$. Put $g_{n,i}=f_i\circ z_n$ and let $M_a$ denote
multiplication by a bounded function $a$ on $H$. By stationarity and the
Markov property,
\begin{equation}\label{eq:fdd-products}
 \EE\prod_{i=0}^k f_i(\xi^n_{t_i})
 =\big\langle 1,
 M_{J_ng_{n,0}}S_n(t_1-t_0)M_{J_ng_{n,1}}
 \cdots S_n(t_k-t_{k-1})M_{J_ng_{n,k}}1\big\rangle_H,
\end{equation}
where $S_n(s)=J_ne^{-sG_n}J_n^*$. Indeed, multiplication by $J_ng_{n,i}$
preserves the range of $J_n$ and represents multiplication by $g_{n,i}$ there.
The tour convergence in Proposition~\ref{prop:inputs} gives
$J_ng_{n,i}\to Jf_i$ uniformly, hence convergence of these multiplication
operators in norm. Equation~\eqref{eq:semigroup-convergence} now passes every
factor in \eqref{eq:fdd-products} to the limit. The resulting expression is
the same product expectation for $B^h$: its semigroup is $e^{-sG}$ and its
initial law $\mu_h$ is stationary because $G1=0$. Products of continuous
functions determine the laws on each finite product of the compact sphere.
Repeated observation times can be combined into a single test function, so
all finite-dimensional distributions converge.
\end{proof}

We shall also need the operator applied to a resolvent approximation. If
$f\in C(\Sph)$, $g_n=f\circ z_n$, and
\begin{equation}\label{eq:resolvent-functions}
 u_{n,\alpha}=\alpha(\alpha+G_n)^{-1}g_n,
 \qquad u_\alpha=\alpha(\alpha+G)^{-1}f,
\end{equation}
then, for every fixed $\alpha>0$,
\begin{equation}\label{eq:resolvent-graph-convergence}
 J_nu_{n,\alpha}\longrightarrow Ju_\alpha,
 \qquad J_nG_nu_{n,\alpha}\longrightarrow JGu_\alpha
 \quad\hbox{in }H.
\end{equation}
The first assertion follows from \eqref{eq:resolvent-convergence} and
$J_ng_n\to Jf$. The second follows from the resolvent equations
$G_nu_{n,\alpha}=\alpha(g_n-u_{n,\alpha})$ and
$Gu_\alpha=\alpha(f-u_\alpha)$.
